\section{Construction of the critical quantum sphere}\label{sec:continuum}

We construct the probability law in Theorem~\ref{thm:continuum} from the
ordinary subcritical unit-area spheres.  The construction precedes, and is
independent of, every discrete-map limit in this paper.  Its three principal
steps are convergence in a maximum-centered cylinder embedding, identification
of area after sampling the normalization marks, and completion of the local
critical metric.  No FK convergence theorem is used in this construction.

\subsection{Local Gaussian inputs and conventions}\label{cont:local-inputs}

All Gaussian free fields have logarithmic singularity
$\log(1/|z-w|)$.  Write
\[
 \mathcal C=\R\times(\R/2\pi\Z),\qquad u=t+i\theta,
 \qquad z=e^u.
\]
The last map identifies the two-point compactification of $\mathcal C$ with
the sphere.  Let $G^\perp$ be the lateral part of a whole-cylinder GFF: its
constant Fourier mode in $\theta$ is removed.  This stationary Gaussian field
can be sampled independently of the average process.  It is also the
nonradial projection of a round mean-zero GFF, transported to the cylinder.
Indeed conformal invariance identifies the Dirichlet spaces, and removal of
the two poles does not change them: logarithmic cutoffs show that a point has
zero Sobolev capacity.  The projection commutes with round rotations about the
poles.  Consequently its scalar transport to the round sphere, with no
background-charge shift, belongs almost surely to $H^{-1}$, since its expected
squared norm is bounded by a constant times
\[
 \sum_{\ell\geq1}\frac{2\ell+1}{[\ell(\ell+1)]^2}<\infty.
\]

We use the following established local Gaussian facts.  For a zero-boundary
GFF $X$ in a disk, the circle-average measures
\[
 \mu_X^\gamma(dz)
   =\lim_{\eps\downarrow0}
      \eps^{\gamma^2/2}e^{\gamma X_\eps(z)}\,d^2z,
       \qquad 0<\gamma<2,
\]
exist and obey the conformal rule with charge $Q_\gamma=2/\gamma+\gamma/2$;
see \cite[Propositions 1.1 and 2.1]{DSKPZ}.  On compact interior sets,
\begin{equation}\label{cont:local-chaos-limits}
 \frac{\mu_X^\gamma}{2(2-\gamma)}
     \xrightarrow[\gamma\uparrow2]{\mathrm{prob.}}\mu_X^{\rm crit},
 \qquad
 (2\log\eps^{-1}-X_\eps)\eps^2e^{2X_\eps}\,d^2z
     \xrightarrow[\eps\downarrow0]{\mathrm{prob.}}\mu_X^{\rm crit}.
\end{equation}
Here convergence is locally weak, and the limit is positive, diffuse, and of
full support.  Moreover the undifferentiated critical approximants
$\eps^2e^{2X_\eps}\,d^2z$ tend locally weakly to zero in probability, with the
usual Seneta--Heyde normalization.  The factor $2$ in
\eqref{cont:local-chaos-limits} is the one in
\cite[Theorem 1.1]{APS}; the remaining critical facts use
\cite{Powell,DRSV,DRSVderivative} with the following circle-average
specialization.
When a result is written in variance rather than radius normalization, the
conversion contributes the deterministic conformal-radius factor.  In the
derivative approximant its additional undifferentiated term has zero limit.
Only interior versions of these Gaussian results are needed below.

\paragraph{The circle-average specialization.}
The inverse-distance mollifier bound printed in \cite[Equation~(1.2)]{Powell}
is not satisfied by uniform measure on a circle, so its use here needs a
direct check. The circle potential identity
\[
 \frac1{2\pi}\int_0^{2\pi}-\log|z-\eps e^{i\theta}|\,d\theta
       =-\log(\max\{|z|,\eps\})
\]
gives, on compact interior sets,
\[
 \Cov(X_\eps(x),X_\delta(y))
 =\log\frac1{|x-y|\vee\eps\vee\delta}+O(1),\qquad
 \E|X_\eps(x)-X_\delta(y)|^2
 \le C\left(\frac{|x-y|}{\eps\wedge\delta}
                   +\left|\log\frac\eps\delta\right|\right)
\]
when $|x-y|\le\eps\wedge\delta$ and $1/2\le\eps/\delta\le2$.
The harmonic part of the Green function preserves these bounds. They supply
the covariance and rescaled H\"older estimates used in
\cite[Lemmas~2.1 and~2.3]{Powell}. For clarity, we give the cutoff-removal argument, using the exact radial
Brownian property of circle averages \cite[Lemma~3.1]{Powell} and the barrier
improvement of \cite[Appendix~A.1]{LacoinCritical}.

Fix a compact interior set and a radius $r_0$ smaller than its distance from
the boundary.  Cover it by $O(e^{2n})$ cells of diameter $e^{-n}$, with centers
$x_i$.  Choose a fixed integer $c$ so that every averaging circle with center
in such a cell and radius in $[e^{-n-1},e^{-n}]$ lies inside the circle centered
at $x_i$ of radius $R=e^{-(n-c)}$.  Subtract $X_R(x_i)$ from all these fine
averages.  The resulting Gaussian process is independent of the centered
coarse radial path down to radius $R$.  Indeed, for $R\le r\le r_0$, the
circle-potential and harmonic mean-value identities give
\[
 \Cov(X_\eps(x),X_r(x_i))=-\log r+g(x,x_i),\qquad
 \Cov(X_R(x_i),X_r(x_i))=-\log r+g(x_i,x_i),
\]
where $g$ is the regular Green kernel.  Their difference is independent of
$r$, and hence has zero covariance with
$X_r(x_i)-X_{r_0}(x_i)$.  This proves independence from the whole centered
coarse path.  The initial field $x\mapsto X_{r_0}(x)$ has finite supremum norm
on the compact set; first restrict $\sup_x|X_{r_0}(x)|\le M$ for a
deterministic $M$.
The displayed increment estimate and Gaussian entropy bounds show that the
supremum $Z_{n,i}$ of the fine residual has uniform subGaussian tails.
In particular its exponential moments of every fixed order are uniformly
bounded.  The coarse Brownian time is $t_n=n+O(1)$.

Convolving the Gaussian endpoint density with this residual bound first
gives, at level $2n+\frac12\log(n+2)+q$, a per-cell probability at most
$C_q e^{-2n}n^{-3/2}$.  Summation over cells and scales gives a rough envelope
with a finite random additive constant.  On the event that this constant is
at most $q$, the coarse Brownian path up to $t_n$ lies below
$2s+L_n$, where $L_n=\frac12\log(n+2)+q+C$.  Retain only this coarse-path
restriction, so that the residual remains independent.  The reflection
formula for Brownian motion killed at this affine barrier bounds its
endpoint density at $u$ by the ordinary Gaussian density times
\[
 \min\left\{1,\frac{C L_n(L_n+2t_n-u)_+}{t_n}\right\}.
\]
At the improved level
$2n-\frac14\log n+2\log\log n+C_q$, the Gaussian density is
$O_q(e^{-2n}(\log n)^{-4})$, and this extra factor, after convolution with
$Z_{n,i}$, is $O_q((\log n)^2/n)$.  Large residuals are controlled by the
same subGaussian bound.  The resulting per-cell probability is therefore
\[
                  \frac{C_q e^{-2n}}{n(\log n)^2}.
\]
It is summable over cells and scales.  Exhausting both deterministic cutoffs
proves that
\[
 \sup_{x\in K,\,0<\eps<\eps_0}
       \bigl(X_\eps(x)-2\log(1/\eps)\bigr)<\infty
       \qquad\text{almost surely}
\]
for every compact interior $K$.  This is the circle-average ceiling needed
for the terminal cutoff removal in \cite[Lemma~3.15]{Powell}.  The fixed
coarse/fine separation is essential to the independence used above.
In the proof of \cite[Section~3]{Powell}, the reference process is now the
circle-average process itself: its comparison parameters satisfy
$\lambda_\eps=1$ and $Y_\eps=\rho_\delta^\eps=0$. The rooted Bessel argument,
GFF Markov decomposition, uniform integrability, and derivative-to-Seneta--Heyde
ratio in Propositions~3.6 and~3.10 therefore apply with these exact parameters.
Circle-average Seneta--Heyde convergence, with the white-noise normalization,
is established in \cite[Theorem~4.2 and its proof]{HRVcriticalDisk}; that proof
explicitly treats the interior Dirichlet GFF. Combining it with the ratio
$\sqrt{2/\pi}$ proves the second limit in
\eqref{cont:local-chaos-limits} with the stated normalization.

\begin{definition}\label{cont:local-ac-definition}
A field on an open planar set is \emph{locally absolutely continuous modulo
constants} if, in a neighborhood of each deterministic point, its restriction
modulo constant functions is absolutely continuous with respect to the
corresponding restriction of a whole-plane GFF modulo constants.  Statements
on an open set use a countable cover.  When several neighborhoods are involved,
we specify joint absolute continuity of their restrictions.
\end{definition}

\begin{lemma}[Transfer of the local area construction]\label{cont:chaos-transfer}
For a field with the local property of
Definition~\ref{cont:local-ac-definition}, the three local chaos limits above
hold on smaller neighborhoods.  The critical limit is a measurable function
of the field.  For every fixed conformal map $f$ from new to old coordinates,
\begin{equation}\label{cont:area-covariance}
 \mu_{h\circ f+2\log|f'|}^{\rm crit}(A)
       =\mu_h^{\rm crit}(f(A))
\end{equation}
on the common coordinate domain.  The local condition itself is preserved by
fixed conformal changes and deterministic smooth additions.
\end{lemma}

\begin{proof}
On a relatively compact patch, whole-plane GFF restrictions modulo constants
and restrictions of a zero-boundary GFF in a larger disk have the same measure
class.  The Markov decomposition supplies the harmonic difference; a smooth
cutoff away from the boundary and the Cameron--Martin theorem give the
equivalence.  The same argument compares different domains and proves the
last assertion using Gaussian conformal invariance.

Normalize a restriction by subtracting its average against a fixed smooth
test function of integral one.  For the reference GFF this changes
subcritical masses by the appropriate exponential of the average.  If $c$ is
a constant and $U_\eps(h)=\eps^2e^{2h_\eps}\,d^2z$, the signed derivative
approximants satisfy
\[
 (2\log\eps^{-1}-(h+c)_\eps)U_\eps(h+c)
 =e^{2c}\bigl[(2\log\eps^{-1}-h_\eps)U_\eps(h)-cU_\eps(h)\bigr].
\]
The last term tends to zero.  Thus all the convergence statements hold for
the normalized representative, transfer by absolute continuity, and remain
true after restoring the possibly random constant.  Restoration can first
be done on bounded-constant events and then exhausted.  Local measurable
versions are obtained by almost-sure subsequences for the reference laws and
uniqueness of the limits; these versions agree on overlaps.

For covariance, apply the subcritical rule to the actual subcritical chaos
of the fixed field, divide by $2(2-\gamma)$, and pass to the limit on both
sides.  Replacing $Q_\gamma\log|f'|$ by $2\log|f'|$ changes the field by a
smooth function tending locally uniformly to zero.  This perturbation is
harmless for positive chaos measures.  More generally, if random continuous
functions $F_\gamma$ converge locally uniformly in probability on the same
probability space, multiplication by $e^{\gamma F_\gamma}$ preserves the
corresponding positive-measure convergence.  The multiplicative identity for
a continuous addition at a fixed subcritical parameter follows directly
from the approximants.  These observations prove
\eqref{cont:area-covariance}.
\end{proof}

\subsection{The maximum-centered subcritical law}\label{cont:shape-section}

Let $W^+$ and $W^-$ be independent standard three-dimensional Brownian motions,
independent of $G^\perp$.  For $a\geq0$ set
\[
 R_a(t)=
 \begin{cases}
   |W^+_t+at e_1|,&t\geq0,\\
   |W^-_{-t}+a(-t)e_1|,&t<0,
 \end{cases}
 \qquad X^a=G^\perp-R_a.
\]
The two definitions agree at zero.  Areas in this subsection are in the flat
cylinder coordinates.  For $\gamma<2$ sufficiently close to two, put
\[
 a=Q_\gamma-\gamma=\frac2\gamma-\frac\gamma2,
 \qquad
 L_\gamma=\frac{\mu_{X^a}^\gamma(\mathcal C)}{2(2-\gamma)}.
\]

\begin{lemma}[Unit-area disintegration at the maximum]\label{cont:shape-law}
Up to horizontal translation, the ordinary DMS unit-area sphere with its two
end marks, shifted by $c_\gamma=\gamma^{-1}\log(2(2-\gamma))$, has cylinder
representative
\begin{equation}\label{eq:original-1}
 X^a-\frac1\gamma\log L_\gamma,
 \qquad
 \frac{d\Prob_\gamma}{d\Prob}
       =\frac{L_\gamma^{2a/\gamma}}{\E L_\gamma^{2a/\gamma}}.
\end{equation}
Adding $c_\gamma$ to the area-one field leaves its normalized area measure
unchanged and multiplies its unnormalized $\gamma$-chaos by
$e^{\gamma c_\gamma}=2(2-\gamma)$.
\end{lemma}

\begin{proof}
The infinite sphere measure underlying
\cite[Definition 4.21(ii)]{DMS} has average process equal to $2/\gamma$ times
a log-Bessel excursion of dimension
$\delta=4-8/\gamma^2\in(0,2)$, on its unit quadratic-variation clock,
independently decorated by $G^\perp$.  The excursion maximum in Bessel units
has intensity proportional to $s^{\delta-3}\,ds$.  Conditional on that
maximum, the two ascents, one read backwards, are independent
$\mathrm{BES}(4-\delta)$ paths from zero to the maximum; this is the Bessel
excursion decomposition at the maximum
\cite[Theorem 1 and the Bessel specialization following it]{PitmanYor}.
For the scaled logarithmic maximum $m$, the substitution
$s=e^{\gamma m/2}$ gives intensity proportional to $e^{-2am}\,dm$.

Place the maximum at horizontal coordinate zero.  On one ascent, the deficit
below $m$, starting at its first visit to a level $K>0$ and ending at zero,
is Brownian motion with drift $-a$, started at $K$ and killed at zero.  This
follows from It\^o's formula and the quadratic-variation clock.  Its reversal
from killing, as $K\to\infty$, has the law of $|W_t+at e_1|$.  We include the
transition-density calculation to fix this identification.

Let $p_b(s,x,y)$ be the killed half-line transition density of Brownian motion
with drift $b$, let $g_b$ be its time integral, and let $q_b(s,x)$ be the
first-hitting density of zero.  At reversed times
$0<s_1<\cdots<s_k$, the subdensity for duration exceeding $s_k$ is
\[
 g_{-a}(K,x_k)
 \left(\prod_{i=2}^k
      p_{-a}(s_i-s_{i-1},x_i,x_{i-1})\right)q_{-a}(s_1,x_1).
\]
The reflected heat kernel with drift gives, for $0<x\leq K$,
\[
 g_{-a}(K,x)=\frac{1-e^{-2ax}}a.
\]
After the drift factors telescope, the limiting density is the zero-drift
$\mathrm{BES}(3)$ density---the same product with $p_0,q_0$ and prefactor
$2x_k$---times
\[
 e^{-a^2s_k/2}\frac{\sinh(ax_k)}{ax_k}.
\]
This is exactly the radial Brownian change of measure for vector drift
$ae_1$.  The duration tends to infinity in probability, and the limiting
density integrates to one.  This identifies the reversed deficit on every
finite horizon.  Both halves remain independent of each other and of $m$;
translation leaves the independent stationary lateral field unchanged.

For a fixed maximum-zero shape with total area $A$, its area at maximum $m$
is $v=e^{\gamma m}A$.  Therefore
\[
 e^{-2am}\,dm
   =\frac1\gamma A^{2a/\gamma}v^{-1-2a/\gamma}\,dv.
\]
Disintegration at $v=1$ tilts the shape by $A^{2a/\gamma}$ and subtracts
$\gamma^{-1}\log A$ from the field.  Replacing $A$ by
$2(2-\gamma)L_\gamma$ changes the tilt only by a deterministic factor;
adding $c_\gamma$ leaves exactly \eqref{eq:original-1}.  This is the
continuous-by-constant-shift version of the area-one disintegration.
Finiteness, positivity, and integrability of its density near two are proved
in the next two lemmas.
\end{proof}

\subsection{Uniform area control on the whole cylinder}\label{cont:tail-section}

Local convergence of chaos does not exclude escape of area toward the two
ends.  We first obtain a small uniform moment, and then sum radial tails.

\begin{lemma}[A uniform positive fractional moment]\label{cont:fractional-moment}
Let $Z_\gamma$ be the area of the inner half of a unit disk for a zero-boundary
GFF, divided by $2(2-\gamma)$.  There exist $s>0$ and $\gamma_0<2$ such that
\[
 \sup_{\gamma_0<\gamma<2}\E Z_\gamma^s<\infty.
\]
After possibly reducing $s$ to lie in $(0,1)$, the normalized areas for
$G^\perp$ on unit cylinder strips have the same uniform bound.
\end{lemma}

\begin{proof}
The normalized masses on slightly larger compact disks are tight by
\eqref{cont:local-chaos-limits}.  Pack $c k^2$ disjoint disks of radius $c'/k$
inside the half-disk.  Their Markov decompositions consist of independent
rescaled zero-boundary fields and an independent collection of harmonic
parts.  With probability at least $3/4$, all those harmonic parts are bounded
below by $-C\log k$ on the corresponding half-disks, uniformly for large $k$.
Indeed their variances on the three-quarter disks are at most
$\log k+O(1)$ by the Green function, or conformal-radius, formula.  Harmonic
interior estimates, the Gaussian exponential estimate and Jensen's inequality
bound each supremum tail; a union bound over $O(k^2)$ disks proves the claim
when $C$ is large.

This event is independent of the inner Dirichlet fields.  On it, scaling
contributes $(c'/k)^{2+\gamma^2/2}$ and hence
\[
 Z_\gamma\geq c_1 k^{-B}\max_{1\leq i\leq c k^2}Z_{\gamma,i}
\]
for a fixed $B$, where the variables on the right are independent copies of
$Z_\gamma$.  Tightness of the parent mass, together with independence of the
harmonic event, yields
\[
 \Prob\!\left(\max_i Z_{\gamma,i}>C_1k^B\right)\leq\tfrac12
\]
uniformly for all sufficiently large $k$, after increasing $C_1$.  Thus
$\Prob(Z_\gamma>C_1k^B)=O(k^{-2})$ uniformly.  Integrating this polynomial
tail proves the first assertion for a sufficiently small positive $s$.

On a compact cylinder patch, realize $G^\perp$ as a zero-boundary field in a
larger coordinate disk plus a continuous random field whose absolute
supremum has every exponential moment.  To do so, add an independent
two-sided average Brownian motion to recover a whole-cylinder field, use
its Markov decomposition modulo constants, and subtract that Brownian
motion.  The remaining harmonic part is smooth on the compact patch and
all normalization constants are Gaussian.  H\"older's inequality, a finite
cover, and the first assertion give a smaller positive moment for the
lateral field.  Stationarity makes this bound uniform over unit strips.
\end{proof}

\begin{lemma}[The cylinder limit]\label{cont:cylinder-limit}
There is a measurable, positive, diffuse, full-support local measure
$m^\perp$ for $G^\perp$ such that, in the Brownian coupling above,
\begin{align*}
 \frac{\mu_{X^a}^\gamma}{2(2-\gamma)}
       &\longrightarrow M_0:=e^{-2R_0}m^\perp,\\
 L_\gamma&\longrightarrow L:=M_0(\mathcal C)\in(0,\infty)
\end{align*}
in probability as $\gamma\uparrow2$.  The first convergence is weak convergence
of finite measures on the compactified cylinder, with no mass at its ends.
There is $s\in(0,1)$ with $\sup_{\gamma\uparrow2}\E L_\gamma^s<\infty$, and
\[
 L_\gamma^{2a/\gamma}\longrightarrow1\quad\hbox{in }L^1.
\]
The same measure convergence holds with $a=0$ on the left.
\end{lemma}

\begin{proof}
On compact cylinder sets, the decomposition used in
Lemma~\ref{cont:fractional-moment}, the positive-measure limit in
\eqref{cont:local-chaos-limits}, and multiplication by the continuous
correction field give convergence to $m^\perp$.  These limits agree on
overlaps, are measurable from $G^\perp$, and are diffuse with full support.
They retain the strip moment bounds by Fatou's lemma.

For each $b>0$, uniformly in $0\leq a\leq1$,
\begin{equation}\label{cont:radial-tail}
 \E\exp\!\left(-b\inf_{[k,k+1]}R_a\right)
       \leq C_b k^{-3/2},\qquad k\geq1,
\end{equation}
and the same estimate holds on the negative half-cylinder.  The density of
$W_k+ak e_1$ is bounded by $Ck^{-3/2}$, independently of its mean.  Integrating
$e^{-b|x|}$ against this bound and then allowing the independent one-unit
increment supremum proves \eqref{cont:radial-tail}; those increments have
uniform exponential moments when $a\leq1$.

The lateral field and the radial processes are independent.  The
$s$-moment of the normalized area on strip $[k,k+1]$ is therefore at most
$Ck^{-3/2}$, using \eqref{cont:radial-tail} with $b=s\gamma$ and taking
$\gamma$ in a fixed interval near two.  The same bound holds for
$e^{-2R_0}m^\perp$.  Since $s<1$, subadditivity shows that the area outside
$|t|\leq K$ has $s$-moment tending uniformly to zero as $K\to\infty$.
On compact sets $R_a\to R_0$ uniformly almost surely.  This proves the weak
convergence, convergence of the totals, and the uniform positive moment.
The limit total is finite by the tail bound and strictly positive by local
full support; its ends have zero mass.

Finally $2a/\gamma\to0$ and $\log L_\gamma\to\log L$ in probability, so the
displayed powers tend to one in probability.  For large $\gamma$ they are
bounded by $1+L_\gamma^{s/2}$, a uniformly integrable family because the
$s$-moments are bounded.  This proves $L^1$ convergence.  The proof of local
convergence and the tail estimate also applies when the left-hand field is
the fixed field $X^0$.
\end{proof}

\subsection{Norm convergence and the three area marks}\label{cont:marked-limit}

\begin{proposition}[The prescribed limiting probability law]\label{cont:limit-law}
The field-measure laws in the continuum approximation of
Theorem~\ref{thm:continuum} converge, independently of the sequence
$\gamma\uparrow2$, in round-background $H^{-1}$ norm and weak measure
topology.  In cylinder coordinates their limit is obtained from
\begin{equation}\label{eq:original-2}
 h_{\mathcal C}=G^\perp-R_0-\tfrac12\log L,
 \qquad \nu_{\mathcal C}=M_0/L,
\end{equation}
by sending the two ends to $0,\infty$ and an additional independent area
sample to $1$, using charge $2$.  Denote this law by $\mathsf P_{\rm crit}$.
It is invariant under resampling three conditionally independent area marks
and normalizing them to $0,1,\infty$.
\end{proposition}

\begin{proof}
In the $z=e^{t+i\theta}$ coordinates, the round-background scalar associated
with \eqref{eq:original-1} is
\[
 G^\perp-R_a-\frac{\log L_\gamma}{\gamma}
                       +Q_\gamma\log\cosh t.
\]
This follows from the flat-cylinder charge rule and
$g_{\rm round}(z)=4/(1+|z|^2)^2$.  The varying radial terms converge in round
$L^2$ in probability.  Indeed the common Brownian bound
$|W^\pm_{|t|}|+O(1+|t|)$ is square-integrable against the exponentially
decaying round-area density at both ends; the random constant converges in
probability by Lemma~\ref{cont:cylinder-limit}.  The lateral field is already
in $H^{-1}$ and is the same in the coupling.  Thus the fields converge in
$H^{-1}$ norm jointly with the finite measures.  The normalized densities
in \eqref{eq:original-1} tend to one in $L^1$, so the conclusion also holds
under the tilted sphere laws.

Under the ordinary area-one DMS sphere law, the two end marks, conditional
on the unmarked surface, are independent area samples
\cite[Proposition A.13]{DMS}.  Appending a third independent area sample and
sending the ends and that sample to $0,\infty,1$, respectively, therefore
gives exactly the three-marked law in the theorem.  The area-sampling kernel
is weakly continuous on probability measures on the compact sphere.  The
limit measure has no atoms, and the third mark avoids the ends, so the
normalizing M\"obius maps converge smoothly as sphere diffeomorphisms.

Composition of scalar distributions with a smoothly convergent family of
sphere diffeomorphisms is continuous in $H^{-1}$: the pullback operators have
locally bounded norms by duality on $H^1$, and the assertion follows first
for smooth distributions and then by approximation.  If $f$ maps new
coordinates to old coordinates, the scalar background correction is
\[
 \frac{Q_\gamma}{2}
 \log\frac{(g_{\rm round}\circ f)|f'|^2}{g_{\rm round}},
\]
which is smooth on the sphere and converges smoothly with $f$ and
$Q_\gamma$.  The continuous mapping theorem now proves the stated joint
limit and its sequence independence.

One may append any finite number of independent area samples before taking
this limit.  At each subcritical parameter, replacing the original three
marks by a fresh triple gives the same law.  The limiting marks remain
distinct almost surely, so passing to the limit in that identity proves the
resampling assertion with charge two.
\end{proof}

\subsection{Local laws and derivative area in sampled coordinates}
\label{cont:local-laws-section}

The measure in \eqref{eq:original-2} has so far been constructed by a
near-critical positive-measure limit.  We next identify its derivative
normalization and justify coordinate changes depending on its own samples.

\begin{lemma}[Absolute continuity on finite slabs]\label{cont:slab-ac}
Translate \eqref{eq:original-2} horizontally so that the maximum of its
average process has an independent smooth location with a strictly positive
density on $\R$.  On every finite slab, the resulting field modulo constants
is absolutely continuous with respect to a whole-cylinder GFF modulo
constants.  The assertion holds jointly on finitely many interior patches.
It also holds for the subcritical representatives in
\eqref{eq:original-1}.
\end{lemma}

\begin{proof}
The translation changes neither the marked embedded surface nor the law of
the independent stationary lateral field.  On a fixed finite horizontal
interval, consider the radial process before subtracting its area
normalization.  If its maximum is outside the interval, the restriction is
absolutely continuous modulo constants with respect to Brownian motion,
because a positive-start $\mathrm{BES}(3)$ path has that property before
approaching zero; time reversal handles the opposite side.  If the maximum
is inside, its location has a density and is therefore absolutely continuous
with respect to the Brownian arcsine maximum-time law.  Conditional on the
location, its two drops are independent negative $\mathrm{BES}(3)$ paths.
The Brownian decomposition at the maximum (apply the minimum decomposition
to $-B$) gives independent meanders for the corresponding drops, and Imhof's relation makes each meander equivalent
to $\mathrm{BES}(3)$ on its prescribed finite time interval; see the sampling
and meander identities in \cite[Corollary~5 and Equations~(5)--(6)]{PitmanMeanders}.  This proves the required
one-dimensional absolute continuity modulo constants.  Adding the independent
lateral field proves the slab assertion, and hence the assertion for joint
interior restrictions.  The normalization $-(\log L)/2$ is only a constant.

For the subcritical shapes, adding vector drift to the Brownian motions is
absolutely continuous on every finite horizon.  The subsequent
absolutely-continuous shape tilt changes no null sets.  The same argument
therefore applies to \eqref{eq:original-1}.
\end{proof}

The local positive-measure limit for the fixed field $X^0$ in
Lemma~\ref{cont:cylinder-limit}, followed by the constant normalization, shows
that $\nu_{\mathcal C}$ is the subcritical-to-critical limit for the fixed
field $h_{\mathcal C}$.  Lemmas~\ref{cont:slab-ac} and
\ref{cont:chaos-transfer} identify it with the derivative circle-average
measure on every finite slab.  These identifications hold jointly with the
independent translation.  Horizontal translations commute exactly with the
circle-average approximants.  On each compact slab their negative mass is
bounded by a finite random ceiling constant times the undifferentiated mass,
which tends to zero; hence their local total variations are tight.
Restricting the translation to a bounded range therefore permits translation
back in the local weak convergence on a fixed larger slab.  Exhausting the
bound gives the same derivative identification for $h_{\mathcal C}$.

\begin{lemma}[Separated area samples]\label{cont:separated-marks}
Take finitely many independent area samples conditional on the cylinder
field.  Restrict their locations to fixed mutually separated compact patches
$U_1,\ldots,U_k$ in a finite cylinder, and let $V$ be an observation patch
separated from them and from the ends.  The joint law of the field near $V$
modulo constants and the sample locations is absolutely continuous with
respect to the local reference GFF law times ordinary Lebesgue measure for
the locations.  Several observation patches may be used simultaneously.
The assertion also holds for the subcritical measures above.
\end{lemma}

\begin{proof}
Consider a set having zero measure for that product reference law.  Its
conditional sampling probability is obtained by integrating its indicator
against the product of the local critical area measures and dividing by a
finite positive power of the total area.  The assertion that its unnormalized
product mass is zero is an event determined by the local field data modulo
constants.  By Lemma~\ref{cont:slab-ac}, it suffices to prove that event has
probability one for a reference cylinder GFF.

For the reference field, condition modulo constants on the complement of
disjoint slightly enlarged domains containing the $U_i$, still separated
from $V$.  The inside fields are independent Dirichlet GFFs.  Their critical
measures almost surely annihilate each deterministic Lebesgue-null set.
Here is a way to see this without an integrability assumption on critical
area.  For a whole-plane GFF modulo constants, translation stationarity
makes the probability of positive mass on a translate of a fixed null set
independent of the translation.  For each realization, local finiteness and
deterministic Fubini show that the mass is zero for almost every translation.
Thus it is zero for any fixed translation.  Local equivalence transfers this
to the interior Dirichlet measures.  In particular their expectations after
normalization to bounded total mass are absolutely continuous with respect
to Lebesgue measure.  Independence gives the same statement for their
product and therefore for any product-Lebesgue-null set of locations.

The conditioned harmonic parts multiply the inside measures by positive
continuous weights; they do not alter this zero-mass conclusion.  The field
on the observation patch is already fixed by the exterior conditioning, and
the product-law null assumption says that the exceptional location section
is Lebesgue-null for almost every such observation field.  Integration proves
the reference assertion, then local absolute continuity proves it for our
field.  Dividing by the total area does not change a zero event.  The same
reasoning, or the usual local first-moment formula, applies subcritically.
\end{proof}

\begin{proposition}[Derivative normalization and area marking]
\label{cont:area-identification}
For $(H,\nu)\sim \mathsf P_{\rm crit}$, put $h=H+\log g_{\rm round}$ in local
conformal coordinates.  On $\widehatC\setminus\{0,1,\infty\}$, its law is
locally absolutely continuous modulo constants and
\[
 \nu(dz)=\mu_h^{\rm crit}(dz)
   =\lim_{\eps\downarrow0}
      (2\log\eps^{-1}-h_\eps(z))\eps^2e^{2h_\eps(z)}\,d^2z
\]
locally weakly in probability.  The measure extends with zero atoms at the
three marks, is diffuse with full support, has total mass one, and is
measurable from $h$.  Its fixed-coordinate covariance and its covariance
under re-marking are the charge-two rule.  The law $\mathsf P_{\rm crit}$ is
nondegenerate.
\end{proposition}

\begin{proof}
We must transfer deterministic-chart statements to maps depending on area
samples.  Cover configurations of old and new marks and observation patches
by countably many location boxes with the separations in
Lemma~\ref{cont:separated-marks} and compact margins for the coordinate maps.
For each deterministic choice of locations in one box, the local Gaussian
law and the chaos statements hold in the observation patch, after subtracting
a test-function mean.  Integrate these statements first against a bounded
location density, then transfer them by the joint absolute continuity of
Lemma~\ref{cont:separated-marks}.  For convergence in probability the same
argument first restricts the Radon--Nikodym density to be bounded and then
exhausts.  Finally restore the random field constant by its scalar
exponential factor; the undifferentiated critical term still vanishes.
This proves both local Gaussian absolute continuity after the random map
and derivative convergence in the common nonsingular region.

For the chart using the two ends and one extra area sample, that region
exhausts the sphere minus $0,1,\infty$.  The local derivative limit there is
the pushforward of $\nu_{\mathcal C}$, by the fixed-field construction and
\eqref{cont:area-covariance}.  The latter probability measure is diffuse and
has full support, and charges none of the three marks.  It is therefore the
unique no-atom extension of the local derivative measures.  Those measures
are local measurable functions of $h$, so the extension is measurable too.
Proposition~\ref{cont:limit-law}'s resampling assertion is consequently an
assertion about samples of exactly this derivative-normalized measure.
The same separated-patch argument proves its covariance under that
re-marking operation.  Finally a deterministic local field restriction has
zero probability under a nontrivial Gaussian field law.  Local absolute
continuity modulo constants therefore precludes a deterministic limiting
field and proves nondegeneracy.  No convergence of subcritical metrics has
been used.
\end{proof}

\subsection{Local critical length metrics}\label{cont:local-metrics-section}

Fix one deterministic normalization of the Ding--Gwynne metric at the
critical LFPP parameter $\xi>0$, characterized by $Q(\xi)=2$.  For a
whole-plane GFF plus a possibly random continuous function, its strong metric
is a measurable function of the field, is a length metric, and is local in
the sense of its internal metrics.  It obeys Weyl scaling, including the
factor $e^{\xi c}$ under an additive constant $c$, and covariance under each
fixed similarity, including rotations
\cite[Definition 1.7, Theorem 1.8, and Proposition 1.9]{DGuniqueness}.
At criticality it is finite and induces the Euclidean topology
\cite[Theorem 1.7]{DGtopology}; the strong metric qualifies for that weak-metric
theorem by \cite[Lemma 1.14]{DGuniqueness}.  Continuous additions preserve
the local topological conclusion by Weyl scaling.  These are the only
metric facts taken as inputs here.  General conformal covariance will be
proved below.

\begin{lemma}[Local versions and pasting]\label{cont:local-metric}
For a field locally absolutely continuous modulo constants, compatible
internal critical metrics can be defined on sufficiently small relatively
compact patches, measurably from the field.  The versions can be chosen
exactly homogeneous under arbitrary constant shifts.  On a connected region
they determine an intrinsic length metric which is finite, induces the
original topology, and has minimizing geodesics between sufficiently near
points.  The construction obeys Weyl scaling for each deterministic smooth
addition and is compatible with each fixed similarity.
\end{lemma}

\begin{proof}
Take as reference $\omega$ a normalized whole-plane GFF plus an independent
Gaussian constant of positive variance.  Its law is equivalent under
deterministic constant shifts and charge-shifted similarities: the
modulo-constant field is similarity invariant, and the remaining constant
has a full-density conditional law.  Locality supplies an internal metric
functional on each patch.  Normalize its input by a fixed smooth
test-function mean, writing the normalized restriction as $\eta$.
Conditional on $\eta$, the reference additive constant has an everywhere
positive density.  Weyl scaling for fixed constants, equivalence under
shifts, and Fubini therefore show that
\[
 c\longmapsto e^{-\xi c}D_{\eta+c}
\]
is essentially constant in $c$, almost surely in $\eta$.  Indeed the
evaluations agree for almost every pair of constants.  Define the local
functional of $\eta$ to be this common essential value, then multiply by
the exponential of the actual mean.  Countable dense distance entries and
continuity give a measurable version.  It agrees with the reference metric
and is exactly constant-homogeneous wherever defined.  Thus it transfers
to arbitrary laws absolutely continuous modulo constants, even when their
actual constants have no density.

For each deterministic smooth addition, the reference Weyl identity is a
local event invariant under a common additive constant.  Smooth additions
preserve the local reference measure class, so the same transfer gives
Weyl scaling for the constructed patch metrics.

Restriction consistency and agreement on overlaps follow from reference
locality and transfer on still smaller patches.  Similarity covariance
transfers in the same way: the transformed reference law uses the same
internal functional by its measure-class equivalence.  We use countable
bases and overlaps within each fixed assertion.

Paste the compatible length data by subdividing paths into patches, or
equivalently by taking infima of finite chain sums of internal distances.
Intermediate chain points may be restricted to a countable dense set,
which proves measurability.  Connectedness gives finite chains.  To exit
a sufficiently small relatively compact neighborhood of a point costs a
strictly positive local distance; hence the pasted distance separates
points and induces the original topology.  For sufficiently near
endpoints, short minimizing sequences cannot exit a fixed compact
neighborhood.  Parametrizing them with bounded speed and applying compactness
gives a local minimizing geodesic.  This argument uses only interior
distances and requires no extension of a patch metric to its boundary.
\end{proof}

\subsection{Conformal covariance at criticality}\label{cont:covariance-section}

We adapt the local linearization and annular comparison strategy of
Gwynne--Miller \cite[Sections~2--3]{GMCovariance}, using the critical
metric's Euclidean topology and the fixed-similarity covariance above.

Fix a univalent conformal map $f$ between coordinate patches.  Let $d$ be the
intrinsic local metric for $h$ and let $\widetilde d$ be the pullback under $f$
of the intrinsic local metric for
\[
 h\circ f^{-1}+2\log|(f^{-1})'|.
\]
Both are continuous-topology length metrics by Lemma~\ref{cont:local-metric}
and Gaussian local absolute continuity.  We will prove equality of their
local lengths.  It suffices first to work with a whole-plane GFF modulo
constants, since the assertions are invariant under constant shifts and
local in the field.

Write $A(a,b)=\{u:a<|u|<b\}$, and set
\[
 A=A(1/2,8),\qquad B=A(0.9,6.5).
\]
For a center $p$ in a fixed compact interior patch and a small radius $r$,
use the old coordinate $u=(z-p)/r$ and the new normalized coordinate
\[
 \psi_{p,r}(u)=\frac{f(p+ru)-f(p)}{rf'(p)}.
\]
This map tends smoothly to the identity, uniformly over the centers on
bounded compact sets.  Similarity covariance removes the same length factor
$r^{2\xi}$ from both metrics; the factor from $f'(p)$ cancels with the charge
term.  After a further common removal of any additive constant, their local
fields are consequently
\begin{equation}\label{cont:annular-fields}
 \eta=h(p+r\,\cdot),\qquad
 \widetilde\eta=\eta\circ\psi_{p,r}^{-1}
                         +2\log|(\psi_{p,r}^{-1})'|.
\end{equation}
For small $r$, the internal metrics needed on $B$ in either coordinate use
only data in a compact subannulus of $A$ in the old coordinates.

\begin{definition}[Successful annulus]\label{cont:successful-annulus}
Fix $\delta_0>0$.  An annulus at $(p,r)$ is successful with constant $C$ if
the following three conclusions hold in the original coordinates.
\begin{enumerate}[label=(\roman*)]
 \item For $x,y$ on the old circles of radii $2r,4r$, respectively,
 $\widetilde d(x,y)$ is at most $1+\delta_0$ times the $d$-length of every
 path joining them within that closed band.
 \item The $d$-distance between any two points of the old $2r$-circle is at
 most $C$ times the infimum of the $d$-lengths crossing the band $[2r,4r]$.
 \item There is a $\widetilde d$-rectifiable loop of winding one about $p$,
 lying in the old open band $(2r,3r)$, whose length is at most $C$ times the
 infimum of the $d$-lengths crossing the band $[r,2r]$.
\end{enumerate}
\end{definition}

\begin{lemma}[One-annulus comparison]\label{cont:single-annulus}
Fix $\delta_0>0$, $M<\infty$, and $q>0$.  Let
$A^+=A(0.75,7.25)$ and require every deterministic smooth addition $u$ to
satisfy $\inf_{a\in\R}\|u-a\|_{C^2(\overline{A^+})}\le M$.
There are $C<\infty$ and
$r_0>0$ such that a zero-boundary GFF on $A$, with any such addition, gives
the sufficient local tests for Definition~\ref{cont:successful-annulus}
probability at least $1-q$, uniformly over the compact set of centers and
$r<r_0$.  Additive constants are unrestricted.
\end{lemma}

\begin{proof}
We first specify the sufficient tests, in rescaled units.  For (i), compare
the two $B$-internal metrics on all pairs of old radii $2$ and $4$, evaluating
the new entries at their $\psi_{p,r}$-images, with factor $1+\delta_0$.
For (ii), bound old $B$-internal entries on the $2$-circle by $C$ times the
infimum of old $B$-internal entries between circles $2$ and $4$.  For (iii),
require a short winding-one loop in the \emph{new} closed ring
$2.25\leq|u|\leq2.75$, with length bounded by $C$ times the corresponding
infimum of old $B$-internal entries between circles $1$ and $2$.
Containment and locality imply the conclusions in the definition: the
internal-distance infima are lower bounds for the constrained crossing
costs, and the new fixed ring lies in the old band $(2,3)$ for small $r$.

These are measurable tests.  In particular, bounded-length winding-loop
existence in a compact ring can be tested using bounded-speed
parametrizations.  For continuous metrics inducing the ring topology,
uniform convergence of metrics and compactness of such paths imply lower
semicontinuity of the required length bound.  Uniform convergence away
from the center preserves winding number.  Thus the loop test is a
measurable condition on the internal metric.

\emph{Uniform comparison of the input laws.}
We claim that the two internal metrics in \eqref{cont:annular-fields},
evaluated at corresponding points, approach one another uniformly in
probability on compact sets of entries in $B$.  First we compare their
field laws.  Choose domains with fixed margins,
\[
 B\Subset W_0\Subset W=A(0.8,7)\Subset A^+\Subset A,
 \qquad W\Subset\psi_{p,r}(A).
\]
For either $D=A$ or $D=\psi_{p,r}(A)$, the Markov decomposition on $W$
writes the relevant restricted field as
\[
 X_W+H_D+u_D,
\]
where $X_W$ is a zero-boundary GFF on $W$, independent of the harmonic
Gaussian field $H_D$.  The smooth term $u_D$ includes the prescribed
addition and, in the transformed coordinate, the logarithmic derivative.
Choose and subtract a constant for which the prescribed addition has
$C^2$ norm at most $M+1$.  Both metrics acquire the same scalar factor,
so all the annular tests are unchanged.

The domains $D$ contain $W$ with uniform margins and lie in one fixed
larger domain.  Differences of their Green functions therefore bound the
variance of $H_D$ uniformly on $W_0$.  Harmonic interior estimates give
uniformly tight smooth norms there.  Choose a smooth cutoff $\chi$ compactly
supported in $W_0$ and equal to one near $\overline B$.  The assumed $C^2$
bound for the additions, together with $\psi_{p,r}\to\mathrm{id}$ smoothly,
implies that the Dirichlet norms of $\chi(H_D+u_D)$ are uniformly tight.
Condition on $H_D$.  On the event that this norm is at most $R$, apply
the Cameron--Martin formula and Cauchy--Schwarz, then integrate over $H_D$.
For every measurable set $E$ of restrictions to $B$, this gives
\[
 \Prob\{(X_W+H_D+u_D)|_B\in E\}
 \le \eps+e^{R^2/2}\Prob\{X_W|_B\in E\}^{1/2},
\]
where $R$ is chosen so that the norm bound fails with probability at most
$\eps$, uniformly in all the parameters.  Thus both input laws are
uniformly absolutely continuous with respect to the same reference
restriction law.

\emph{Continuity in probability of the metric comparison.}
Under their original coupling the two fields in
\eqref{cont:annular-fields} approach one another in a local negative Sobolev
topology, uniformly over the stated centers and additions.  Apply Lusin's
theorem to the measurable map from the restriction
to its $B$-internal metric, viewed as a continuous distance function on each
compact set of entries.  On a compact set of arbitrarily large reference
probability this map is uniformly continuous.  The preceding uniform
absolute continuity controls the probability that either input leaves
that set, and closeness of the coupled inputs then controls the two
metrics.  Their tight continuity moduli also absorb the endpoint motions
by $\psi_{p,r}$.  This proves the claim; it uses measurability of the metric
functional and comparison of the input laws.

\emph{The three geometric tests.}
Internal-distance infima between two disjoint compact circles are strictly
positive, and compact suprema are finite.  Consequently the first two tests
have probability arbitrarily near one when $C$ is large and then $r$ is
small.  The least length of an admissible winding loop is tight by the
same uniform absolute continuity.  For each continuous-topology length
metric, some finite-length winding loop exists: approximate a middle
Euclidean circle by joining sufficiently many successive nearby points
using short local paths which remain close to the circle.  This proves
the third test, uniformly under the stated bounds on the additions.
\end{proof}

\begin{lemma}[Successful annuli at every point]\label{cont:all-point-annuli}
For a fixed conformal map and fixed $\delta_0>0$, almost surely every point
of a compact interior patch has arbitrarily small successful annuli with
one common constant $C$.  Each annulus center is within $r/2$ of its point.
The assertion may be imposed simultaneously on a countable compact
exhaustion and a sequence $\delta_0\downarrow0$.
\end{lemma}

\begin{proof}
Fix $\lambda>20$ and, for a deterministic center $p$, use
$r_j=\lambda^{-j}$, $m\leq j\leq2m$.  The annuli $p+r_jA$ are disjoint.
Decompose the whole-plane field into their independent Dirichlet pieces
and the independent projection harmonic on all of them, modulo constants.
Let $H_j$ denote the rescaled harmonic restriction, normalized by subtracting
a smooth average on $A''=A(0.6,7.8)$.  For all sufficiently large $N$,
\begin{equation}\label{cont:harmonic-boost}
 \Prob\!\left(\sum_{j=m}^{2m}\|H_j\|_{L^2(A'')}^2>Nm\right)
       \leq e^{-cNm}.
\end{equation}
The constant is uniform in $p,m$ on the compact patch.  To verify the bound,
the diagonal covariance kernels are uniformly bounded after the Dirichlet
log singularity and the common average have been removed.  For distinct
annuli, the cross-covariance is the whole-plane logarithmic kernel
double-centered by the averages; it is
$O(\lambda^{-|j-k|})$.  The covariance on the direct-sum $L^2$ space
therefore has bounded operator norm and trace $O(m)$.  Diagonalization and
a small fixed exponential moment of its squared Gaussian norm give
\eqref{cont:harmonic-boost}.

Outside this exception, at least half the indices have bounded harmonic
$C^2$ norm on $\overline{A^+}=\overline{A(0.75,7.25)}$, with a bound
depending only on $N$.  Harmonic interior estimates give this bound from
the $L^2(A^{\prime\prime})$ bound, since $\overline{A^+}\subset A^{\prime\prime}$.
Thus the exact norm and fixed margins required by
Lemma~\ref{cont:single-annulus} hold.
Conditional on the common harmonic projection, the corresponding tests
use independent Dirichlet pieces.  For those indices their failure
probabilities are at most any prescribed $q>0$, once $C$ and then the
minimum scale have been chosen.  Hence the probability of no success at
the center is at most
\[
 e^{-cNm}+q^{m/2}.
\]
Choose $N$ first, then $q$, so that this is bounded by $e^{-c_*m}$ with
$c_*>5\log\lambda$, increasing the lower cutoff on $m$ if necessary.

At level $m$, a lattice of spacing $\lambda^{-2m}/10$ has
$O(\lambda^{4m})$ centers near the compact patch.  The union bound and
Borel--Cantelli show that, almost surely, every such center has a success
for all sufficiently large $m$.  Every point is within
$\lambda^{-2m}/10<r_j/2$ of a lattice center and hence has a successful
annulus enclosing it at a scale tending to zero.  Only countably many
deterministic centers and radii were used for similarity covariance.
Finally exhaust compact patches and intersect over the prescribed
countable sequence of errors.  No quantitative rate in the Lusin argument
was required: its failure probability was chosen before the minimum scale.
\end{proof}

\begin{proposition}[Fixed conformal covariance]\label{cont:metric-covariance}
For every fixed conformal coordinate map, the local critical metrics of
Lemma~\ref{cont:local-metric} obey the charge-two coordinate-change rule.
Equivalently, the metrics $d$ and $\widetilde d$ above have the same local
lengths and hence the same intrinsic distances on common domains.
\end{proposition}

\begin{proof}
First work with the reference field and a fixed $\delta_0>0$ on the
probability-one event in Lemma~\ref{cont:all-point-annuli}.  We give two
steps, since no Euclidean H\"older estimate is available or needed.

\emph{A finite Lipschitz bound along geodesics.}
Let $P:[0,\ell]\to U$ be any nonconstant minimizing $d$-geodesic in the patch,
parametrized by $d$-length.  At each interior time choose a successful small
loop surrounding $P(t)$, with both endpoint images outside its enclosed
region.  Let the nearest hits on this loop before and after $t$ delimit
an open time interval.  Throughout that interval $P$ remains in a bounded
complementary component with nonzero winding.  The loop's $\widetilde d$
length is at most $C$ times the interval's duration: the geodesic crosses
the corresponding band $[r,2r]$ there.

For $0<\eps<\ell/2$, finitely many such intervals cover
$[\eps,\ell-\eps]$, with all loop diameters as small as desired.  Prune to
a minimal interval cover.  Its intervals form an ordered chain without
containments and have multiplicity at most two.  Consecutive loops
intersect.  Indeed their hit intervals interlace, so a point of each loop
lies in a bounded complementary component of nonzero winding of the other.
If the connected loops were disjoint, each entire loop would lie in such
a component of the other.  This is impossible, for example by taking their
extreme horizontal coordinates.  The argument does not require simple
loops.  Concatenating along the intersecting chain costs at most $2C\ell$.
Let the loop mesh tend to zero and then $\eps\downarrow0$.  Continuity gives
\[
 \widetilde d(P(0),P(\ell))\leq2C\ell.
\]
The argument applies to every subsegment.  Thus
$F(t)=\widetilde d(P(0),P(t))$ is Lipschitz on $[0,\ell]$.

\emph{The sharp bound at differentiability times.}
At an interior differentiability time $t$ of $F$, take successful annuli
shrinking to $P(t)$.  Let $u<t<v$ be the nearest visits on either side to
old radius $4r$, and let $k,s$ be the first and last visits to radius $2r$
between them.  All four times tend to $t$, by simplicity and compactness
of a minimizing geodesic.  Its minimality and test (ii) imply
\[
 s-k=d(P(k),P(s))\leq C(k-u).
\]
Both $[u,k]$ and $[s,v]$ are traversals eligible for test (i), and
\[
 v-u\leq(C+2)\max\{k-u,v-s\}.
\]
One eligible traversal therefore occupies a fixed positive fraction of
the whole shrinking interval.  Test (i) bounds the absolute change of $F$
on it by $1+\delta_0$ times its duration.  Differentiability at $t$ has
error $o(v-u)$, which remains negligible after division by that traversal
duration.  Hence $|F'(t)|\leq1+\delta_0$.  Integrating the Lipschitz function
gives
\[
 \widetilde d(P(0),P(\ell))\leq(1+\delta_0)\ell.
\]
Local geodesic accessibility proves this bound for all sufficiently near
points.  Pass through the countable sequence $\delta_0\downarrow0$.
Apply the same argument in the opposite coordinates, whose field is
locally absolutely continuous modulo constants, to obtain the reverse
inequality.  The germs agree, and therefore all intrinsic lengths agree.
Local absolute continuity transfers the result from the reference field
to the stated class of fields.
\end{proof}

\begin{corollary}[Short surrounding detours and random charts]
\label{cont:detours}
Almost surely, every regular point has arbitrarily small surrounding
rectifiable loops whose metric lengths tend to zero.  The conformal
covariance of Proposition~\ref{cont:metric-covariance} also holds on common
nonsingular regions for the random maps obtained from finitely many area
samples in the preceding construction.
\end{corollary}

\begin{proof}
Use the third successful-annulus test with $f=\id$.  Fix a local finite
minimizing segment ending at the point under consideration.  The cost of
crossing the shrinking displaced bands is bounded by the length of a tail
of that segment and hence tends to zero.  Lemma~\ref{cont:all-point-annuli}
holds simultaneously at every point, so this argument gives the same
all-point quantifier for the loops.  It transfers by local absolute
continuity.

For random maps, use Lemma~\ref{cont:separated-marks} on countably many
configuration boxes and observation patches avoiding all relevant marks.
Internal metric comparisons on smaller neighborhoods are measurable local
tests, including as the M\"obius parameters vary; one can use fixed countable
coordinate bases.  For each deterministic parameter the proposition holds
almost surely.  Integration against the product reference law and then
joint absolute continuity transfer it to the sampled parameters.  This is
why fixed-map almost-sure covariance alone is not the argument for a
mark-dependent coordinate change.
\end{proof}

\subsection{Completion at the marked points}\label{cont:completion-section}

\begin{proposition}[The metric sphere]\label{cont:sphere-completion}
The intrinsic critical metric on
$\widehatC\setminus\{0,1,\infty\}$ for $\mathsf P_{\rm crit}$ completes to a finite
continuous metric $D_h$ on the sphere.  Its completion adds exactly the three
marked points and induces the usual sphere topology.  The extended metric
is measurable from the original field alone and obeys the charge-two
coordinate rule, including the re-marking operation.
\end{proposition}

\begin{proof}
Couple the original three-marked sphere with three fresh independent area
samples, and express the same abstract sphere in the coordinate defined by
the fresh triple.  Diffuseness makes the two triples disjoint almost surely.
By Proposition~\ref{cont:limit-law}, the field in each coordinate has law
$\mathsf P_{\rm crit}$.  Pull both local length structures back to the
common abstract sphere.  Each is defined away from its own exceptional
triple, and their lengths agree away from the union of the two triples by
Corollary~\ref{cont:detours}.  Its random-map assertion applies using cylinder
observation patches separated from the ends and all sampled locations.

The two regular regions cover the abstract sphere.  Their compatible
local length structures therefore paste as in
Lemma~\ref{cont:local-metric}, giving a finite length metric inducing the
sphere topology.  This metric space is compact.  Every point of either
exceptional triple has arbitrarily short surrounding detours because it
is regular for the other coordinate.  It remains to show that removing
one triple does not change intrinsic distances between the remaining
points.

For a path whose endpoints avoid one triple, choose short separating loops in
arbitrarily small disjoint neighborhoods of its three points, avoiding the
path endpoints.  Replace the path between its first and last encounter
with each loop by an arc of that loop.  The increase in length is at most
the loop length, and no new visit to the exceptional triple is introduced.
Sending these costs to zero shows that intrinsic distances avoiding the
triple equal the pasted distances there.

The original punctured metric is thus exactly the restriction of this
compact metric sphere.  Its completion adds the three distinct missing
points and no others.  It is determined measurably by the original field:
construct its punctured local metrics using countable entries and take
their unique continuous extension.  The auxiliary chart proves the
extension property and contributes no metric randomness.  Covariance on
the common regular regions extends uniquely over the finite exceptional
sets.  The same applies to the already identified area, proving the
asserted fixed-chart and re-marking rules.
\end{proof}

\begin{remark}[Simultaneous versions on a marked M\"obius orbit]
\label{cont:orbit-versions}
If simultaneous coordinate versions are desired, track the three removed
points and equivariantize the measurable area and metric functionals.
Choose a probability equivalent to Haar measure on the M\"obius group,
evaluate the functionals after each coordinate change, and pull them back.
Fixed-map covariance and Fubini show that, almost surely for the present
law, these outputs have a common essential value agreeing with the
original output.  Take that value.  Invariance of the Haar measure class
under translation transports its existence and value pathwise under every
further M\"obius change, giving a simultaneous version on the marked orbit.

Joint measurability can be expressed using area on rational chart disks and
the internal metrics pasted from disks whose closures avoid the tracked
points.  Continuous extension is checked on countable dense sets, with
dummy outputs when a construction fails.  The needed modulo-constant
absolute continuity on such disks also follows directly from
Lemma~\ref{cont:separated-marks}: in each location box its preimage lies in
a slightly enlarged observation domain separated from the marks.  In the
normalized embeddings the tracked triple is deterministic.
\end{remark}

\subsection{Signed approximants on deterministic compact sets}
\label{cont:compact-restrictions-section}

For completeness we record a version of derivative convergence which allows
restriction to a deterministic compact set without a prior boundary
continuity assumption.

\begin{lemma}[Control of the negative part]\label{cont:compact-restrictions}
On the local absolutely continuous patches above, the negative parts of the
signed derivative approximants tend to zero in probability, and their total
variations on compact patches are tight.  For every deterministic bounded
Borel test function with compact support in a patch, its integral against
the approximants converges in probability to its integral against
$\mu_h^{\rm crit}$.  In particular one has weak
convergence in probability after restriction to any deterministic compact
subset of a patch.
\end{lemma}

\begin{proof}
Use a centered flat-torus GFF as reference.  It is translation invariant
and has the same interior measure class modulo constants.  Set
$T=\log\eps^{-1}$.  Circle-average variance is at most $T+O(1)$, and the
variance of an increment between centers at distance $b\leq\eps$ is
$O(b/\eps)$.  This follows from the logarithmic kernel plus its smooth
remainder: one circle average already gives a Lipschitz bound
$O(\eps^{-1})$ for the logarithmic part.

Gaussian union bounds on a grid of spacing $\eps T^{-2}$ give an upper
bound $2T+O(\log T)$ for the grid maximum with probability tending to one.
Use dyadic grid refinements, with increment thresholds $2^{-j/4}$ at
refinement $j$.  Their variances are $O(T^{-2}2^{-j})$; the Gaussian tails,
including the number of grid edges, are summable in $j$.  Continuity then
gives
\[
 \max h_\eps\leq2T+O(\log T)
 \quad\hbox{with probability tending to one}.
\]
Seneta--Heyde convergence gives undifferentiated total critical mass
$O_{\Prob}(T^{-1/2})$.  The negative derivative mass is at most this mass
times $(\max h_\eps-2T)_+$, hence is
$O_{\Prob}(\log T/\sqrt T)=o_{\Prob}(1)$.  Convergence of the signed total
mass now implies tightness of the total variation.

Divide the variation measure by one plus its total mass.  The resulting
bounded random measure is stationary, so its expectation is a constant
multiple of Lebesgue measure with uniformly bounded density.  Together with
variation tightness this gives uniformly small mass in probability on
deterministic sets of small Lebesgue measure.  The limiting measure has the
same property after the analogous normalization.  Approximate a bounded
Borel test function in Lebesgue measure by bounded continuous functions;
these small-set bounds and tightness control the error uniformly, and local
weak convergence handles the continuous approximations.  This proves the
Borel assertion.  Applying it to a compact-set indicator times a
continuous test function, then using tight variation, proves weak
convergence on each deterministic compact restriction.

Normalize by a local test-function mean to transfer through absolute
continuity, then restore the actual constant using the identity in the
proof of Lemma~\ref{cont:chaos-transfer}.  The undifferentiated term still
vanishes.  A finite cover of a compact patch completes the transfer.
\end{proof}

\begin{proof}[Proof of Theorem~\ref{thm:continuum}]
Proposition~\ref{cont:limit-law} constructs the unique limit of the prescribed
subcritical probability laws in the required norm field and weak measure
topologies.  Proposition~\ref{cont:area-identification} supplies its
nondegeneracy, derivative normalization, full support, diffuseness,
unit total mass, measurability, and area-mark resampling rule.
Lemma~\ref{cont:compact-restrictions} gives the stronger compact-restriction
interpretation if used.  Proposition~\ref{cont:sphere-completion} constructs
the required field-measurable finite critical metric and identifies its
completion and coordinate rules.

Finally take the product of the entire marked $(H,\nu)$ law with the fixed
whole-plane nested $\CLE_4$ probability.  Its M\"obius invariance is
\cite[Theorem 1]{KWnested}, whose parameter range includes four.  This
defines the specified independent decorated continuum law.  No discrete
subsequence or assumption about convergence of subcritical metrics enters
any part of this construction.
\end{proof}
